\documentclass[11pt]{amsart}
\usepackage[margin=1in]{geometry}
\usepackage{amsmath,amssymb,amsthm,mathtools}
\usepackage[T1]{fontenc}
\usepackage{lmodern}
\usepackage{microtype}
\usepackage{enumitem}
\usepackage[hidelinks]{hyperref}

\newtheorem{theorem}{Theorem}[section]
\newtheorem{lemma}[theorem]{Lemma}
\newtheorem{proposition}[theorem]{Proposition}
\newtheorem{corollary}[theorem]{Corollary}
\theoremstyle{definition}
\newtheorem{definition}[theorem]{Definition}
\newtheorem{example}[theorem]{Example}
\theoremstyle{remark}
\newtheorem{remark}[theorem]{Remark}

\newcommand{\R}{\mathbb{R}}
\newcommand{\N}{\mathbb{N}}
\newcommand{\Z}{\mathbb{Z}}

\title[The Canonical Reciprocal Cost]{The Canonical Reciprocal Cost:\\
Exact Multi-Bond Minimization, Frustration Lower Bounds,\\
and Simultaneous-vs-Sequential Descent}

\author{Jonathan Washburn}
\address{Austin, Texas, USA}
\email{washburn.jonathan@gmail.com}
\date{February 2026}

\begin{document}
\begin{abstract}
We study the reciprocal cost
\[
J(x):=\frac12\!\left(x+\frac1x\right)-1,\qquad x>0.
\]
After proving its elementary structure (non-negativity, reciprocal
symmetry, strict convexity, unit curvature at the minimum, and the
log-domain identity $J(e^\varepsilon)=\cosh\varepsilon-1$), we derive an
exact closed-form formula for the multi-bond minimizer of
$\mathcal C(v)=\sum_{i=1}^N J(v/n_i)$:
\[
v^\star=\sqrt{\frac{\sum_i n_i}{\sum_i n_i^{-1}}}
=\sqrt{\operatorname{AM}\cdot\operatorname{HM}}.
\]
The optimizer is thus the geometric mean of the arithmetic and harmonic
means of the neighbor data (and equals the ordinary geometric mean when $N=2$).
We prove a strict frustration inequality: for incompatible neighbors
$n_1\neq n_2$, the exact two-neighbor residual is
\[
\min_{v>0}\bigl(J(v/n_1)+J(v/n_2)\bigr)
=\sqrt{\frac{n_1}{n_2}}+\sqrt{\frac{n_2}{n_1}}-2>0.
\]
For dynamics, we show that continuous-time simultaneous gradient flow
converges to $v^\star$ (via a Lyapunov argument), while cyclic sequential
single-bond projection is periodic with Ces\`aro mean equal to the
arithmetic mean; for distinct neighbors these limits differ by strict AM-GM.
Finally, we record the multiplicative d'Alembert functional identity for~$J$,
characterize $J$ as the unique even solution satisfying this identity with
unit curvature, and derive that the golden ratio $\varphi=(1+\sqrt5)/2$
yields the minimal nonzero penalty scale $\ln\varphi\approx 0.4812$.
\end{abstract}

\subjclass[2020]{Primary 26A51; Secondary 39B22, 52A41, 26D15}
\keywords{Convex reciprocal cost, strict frustration, arithmetic-harmonic mean,
d'Alembert functional equation, gradient flow, AM-GM inequality, golden ratio}
\maketitle

%=======================================================================
\section{Introduction}\label{sec:intro}
%=======================================================================

Define
\begin{equation}\label{eq:J-def}
J:(0,\infty)\to\R,\qquad
J(x):=\frac12\!\left(x+\frac1x\right)-1.
\end{equation}
This function has appeared in various guises across mathematics:
it is the pullback of $\cosh-1$ under the exponential map (Proposition~\ref{prop:cosh}),
the unique normalized solution of the multiplicative d'Alembert functional
equation (Theorem~\ref{thm:RCL} and Remark~\ref{rem:uniqueness}),
and the natural reciprocal-invariant cost on the multiplicative group
$(0,\infty)$ with the Haar measure $dx/x$; see Acz\'el--Dhombres~\cite{AczelDhombres}
for the functional-equation perspective and
Hardy--Littlewood--P\'olya~\cite{HardyLittlewoodPolya} for the
inequality-theoretic background.

\medskip

This paper has four goals, each building on the previous.
\begin{enumerate}[label=(\Roman*),nosep]
\item \textbf{Core calculus} (\S\ref{sec:core}):
  establish $J\ge0$, $J(x)=J(1/x)$, strict convexity, the cosh identity,
  and unit curvature $J''(1)=1$.
\item \textbf{Exact minimization} (\S\ref{sec:minimization}):
  for positive neighbor data $n_1,\dots,n_N$, compute the exact
  minimizer of $\mathcal C(v)=\sum_i J(v/n_i)$ in closed form.
\item \textbf{Frustration inequalities} (\S\ref{sec:frustration}):
  prove that incompatible neighbors force a strictly positive residual
  cost, with an explicit formula for the minimum.
\item \textbf{Descent comparison} (\S\ref{sec:descent}):
  show that simultaneous descent converges to the $J$-optimal point
  while sequential descent converges (in time-average) to the strictly
  suboptimal arithmetic mean.
\end{enumerate}
All proofs are elementary and self-contained; the only prerequisites are
single-variable calculus and the Cauchy--Schwarz inequality.

%=======================================================================
\section{Core identities and convexity}\label{sec:core}
%=======================================================================

\begin{lemma}[Equivalent square form]\label{lem:square}
For every $x>0$,
\begin{equation}\label{eq:square}
J(x)=\frac{(x-1)^2}{2x}.
\end{equation}
\end{lemma}
\begin{proof}
Expand the numerator and divide term by term:
\[
\frac{(x-1)^2}{2x}
=\frac{x^2-2x+1}{2x}
=\frac{x}{2}-1+\frac{1}{2x}
=\frac12\!\left(x+\frac1x\right)-1
=J(x).  \qedhere
\]
\end{proof}

\begin{proposition}[Basic structure]\label{prop:basic}
For all $x>0$:
\begin{enumerate}[label=\textup{(\alph*)},nosep]
\item\label{it:unit} $J(1)=0$.
\item\label{it:nn} $J(x)\ge0$, with equality if and only if $x=1$.
\item\label{it:sym} $J(x)=J(1/x)$ \textup{(reciprocal symmetry)}.
\item\label{it:deriv} $J'(x)=\frac12(1-x^{-2})$, $J''(x)=x^{-3}>0$, $J''(1)=1$.
\end{enumerate}
In particular, $J$ is strictly convex on $(0,\infty)$ with unique global
minimizer at $x=1$.
\end{proposition}
\begin{proof}
\ref{it:unit}: $J(1)=\frac12(1+1)-1=0$.

\ref{it:nn}: By Lemma~\ref{lem:square}, $J(x)=(x-1)^2/(2x)$.
The numerator $(x-1)^2$ is a perfect square, hence $\ge0$.
The denominator $2x$ is strictly positive for $x>0$.
Therefore $J(x)\ge0$.  Equality forces $(x-1)^2=0$, i.e.\ $x=1$.

\ref{it:sym}: Replacing $x$ by $1/x$ in the definition:
\[
J(1/x)=\frac12\!\left(\frac1x+x\right)-1=J(x).
\]

\ref{it:deriv}: Differentiating~\eqref{eq:J-def}:
\[
J'(x)=\frac12\!\left(1-x^{-2}\right),\qquad
J''(x)=\frac12\cdot 2x^{-3}=x^{-3}.
\]
Since $x^{-3}>0$ for all $x>0$, the function is strictly convex.
At $x=1$: $J''(1)=1$.  (The value $J''(1)=1$ means $J$ has unit curvature
at its minimum; this is the normalization that makes $J$ canonical among
even convex functions on the multiplicative line.)
\end{proof}

\begin{proposition}[Log-domain representation]\label{prop:cosh}
For every $\varepsilon\in\R$,
\begin{equation}\label{eq:cosh}
J(e^\varepsilon)=\cosh\varepsilon-1.
\end{equation}
Consequently, $J(e^\varepsilon)=\frac{\varepsilon^2}{2}+\frac{\varepsilon^4}{24}
+\frac{\varepsilon^6}{720}+\cdots$ \textup{(}convergent for all $\varepsilon$\textup{)}.
\end{proposition}
\begin{proof}
Substitute $x=e^\varepsilon$ into the definition:
\[
J(e^\varepsilon)
=\frac12\!\left(e^\varepsilon+e^{-\varepsilon}\right)-1
=\cosh\varepsilon-1.
\]
The Taylor series $\cosh\varepsilon=\sum_{k=0}^\infty\varepsilon^{2k}/(2k)!$
converges absolutely for all $\varepsilon\in\R$, giving the stated expansion.
\end{proof}

\begin{remark}[Geometric interpretation]\label{rem:geom}
Equation~\eqref{eq:cosh} shows that in the \emph{log domain}
(writing $x=e^\varepsilon$, so that the multiplicative group $(0,\infty)$
becomes the additive line $\R$), the cost $J$ is simply a shifted
hyperbolic cosine.  This is the even, convex, unit-curvature potential
on $(\R,+)$, exactly as $J$ is the even, convex, unit-curvature potential
on $((0,\infty),\times)$.  The translation between additive and
multiplicative pictures is a recurring theme.
\end{remark}

\begin{corollary}[Small-strain approximation]\label{cor:small}
For $|\varepsilon|\le1$,
\[
J(1+\varepsilon)=\frac{\varepsilon^2}{2(1+\varepsilon)}
=\frac{\varepsilon^2}{2}-\frac{\varepsilon^3}{2}+O(\varepsilon^4).
\]
\end{corollary}
\begin{proof}
By Lemma~\ref{lem:square}, $J(1+\varepsilon)=\varepsilon^2/(2(1+\varepsilon))$.
Expand $(1+\varepsilon)^{-1}=1-\varepsilon+\varepsilon^2-\cdots$ for $|\varepsilon|<1$.
\end{proof}

\begin{corollary}[Ratio-zero criterion]\label{cor:ratio-zero}
For $v,n>0$,
\[
J(v/n)=0 \iff v=n.
\]
\end{corollary}
\begin{proof}
$J(v/n)=0\iff v/n=1$ by Proposition~\ref{prop:basic}\ref{it:nn}.
\end{proof}

%=======================================================================
\section{Exact multi-bond minimization}\label{sec:minimization}
%=======================================================================

Given positive ``neighbor values'' $n_1,\dots,n_N$ ($N\ge1$), we seek the
value $v>0$ that minimizes the total bond cost.  The key insight is that
$J$'s linearity in $v$ and $1/v$ (after expanding each term) yields a
clean closed form.

\begin{definition}[Total bond cost]\label{def:total}
For $n_1,\dots,n_N>0$ and $v>0$, define
\[
\mathcal C(v):=\sum_{i=1}^N J(v/n_i).
\]
\end{definition}

\begin{theorem}[Closed form and unique minimizer]\label{thm:minimizer}
Set
\[
A:=\sum_{i=1}^N\frac1{n_i},\qquad
B:=\sum_{i=1}^N n_i.
\]
Then:
\begin{enumerate}[label=\textup{(\roman*)},nosep]
\item\label{it:closed}
$\displaystyle\mathcal C(v)=\frac12\!\left(Av+\frac{B}{v}\right)-N.$
\item\label{it:convex}
$\mathcal C$ is strictly convex on $(0,\infty)$ with
$\mathcal C(v)\to+\infty$ as $v\to 0^+$ or $v\to+\infty$.
\item\label{it:min}
The unique minimizer is
\begin{equation}\label{eq:vstar}
v^\star=\sqrt{\frac{B}{A}}
=\sqrt{\frac{\sum_i n_i}{\sum_i n_i^{-1}}}.
\end{equation}
\end{enumerate}
\end{theorem}
\begin{proof}
\ref{it:closed}: Expand each summand using the definition of $J$:
\[
J(v/n_i)=\frac12\!\left(\frac{v}{n_i}+\frac{n_i}{v}\right)-1.
\]
Summing over $i=1,\dots,N$:
\[
\mathcal C(v)
=\frac{v}{2}\sum_i\frac{1}{n_i}+\frac{1}{2v}\sum_i n_i-N
=\frac12\!\left(Av+\frac{B}{v}\right)-N.
\]

\ref{it:convex}: Differentiate the closed form:
\[
\mathcal C'(v)=\frac12\!\left(A-\frac{B}{v^2}\right),\qquad
\mathcal C''(v)=\frac{B}{v^3}.
\]
Since $B=\sum n_i>0$ and $v>0$, we have $\mathcal C''(v)>0$ everywhere,
giving strict convexity.  As $v\to0^+$, the term $B/(2v)\to+\infty$;
as $v\to+\infty$, the term $Av/2\to+\infty$.  So $\mathcal C$ is coercive.

\ref{it:min}: A strictly convex coercive function on $(0,\infty)$ has
exactly one critical point, which is the global minimum.
Setting $\mathcal C'(v)=0$:
\[
A=\frac{B}{v^2}\implies v^2=\frac{B}{A}\implies v=\sqrt{B/A}.
\qedhere
\]
\end{proof}

\begin{corollary}[Mean interpretation]\label{cor:means}
With $\operatorname{AM}:=\frac1N\sum_i n_i$ and
$\operatorname{HM}:=N/\!\sum_i n_i^{-1}$,
\[
v^\star=\sqrt{\operatorname{AM}\cdot\operatorname{HM}}.
\]
In particular, for $N=2$: $v^\star=\sqrt{n_1n_2}$ (the geometric mean).
\end{corollary}
\begin{proof}
From~\eqref{eq:vstar}:
\[
v^\star=\sqrt{\frac{B}{A}}
=\sqrt{\frac{N\cdot\operatorname{AM}}{N/\operatorname{HM}}}
=\sqrt{\operatorname{AM}\cdot\operatorname{HM}}.
\]
For $N=2$: $B/A=(n_1+n_2)/(n_1^{-1}+n_2^{-1})
=(n_1+n_2)\cdot n_1n_2/(n_1+n_2)=n_1n_2$, so $v^\star=\sqrt{n_1n_2}$.
\end{proof}

\begin{corollary}[Minimal cost value]\label{cor:minvalue}
\[
\min_{v>0}\mathcal C(v)
=\mathcal C(v^\star)
=\sqrt{AB}-N\ge0.
\]
Equality $\sqrt{AB}=N$ holds if and only if $n_1=\cdots=n_N$.
\end{corollary}
\begin{proof}
Substitute $v^\star=\sqrt{B/A}$ into the closed form:
\[
\mathcal C(v^\star)
=\frac12\!\left(A\sqrt{B/A}+\frac{B}{\sqrt{B/A}}\right)-N
=\frac12\!\left(\sqrt{AB}+\sqrt{AB}\right)-N
=\sqrt{AB}-N.
\]
By the Cauchy--Schwarz inequality applied to the vectors
$(n_1,\dots,n_N)$ and $(1/n_1,\dots,1/n_N)$ in $\R^N$:
\[
AB=\Bigl(\sum_i n_i\Bigr)\!\Bigl(\sum_i \frac1{n_i}\Bigr)
\ge\Bigl(\sum_i\sqrt{n_i\cdot\frac1{n_i}}\Bigr)^2=N^2.
\]
Hence $\sqrt{AB}\ge N$, giving $\mathcal C(v^\star)\ge0$.
Equality in Cauchy--Schwarz holds iff $(n_i)$ and $(1/n_i)$ are
proportional, i.e.\ iff all $n_i$ are equal.
\end{proof}

\begin{example}\label{ex:two-bond}
Let $n_1=1$, $n_2=4$.  Then $A=1+1/4=5/4$, $B=1+4=5$,
$v^\star=\sqrt{5/(5/4)}=\sqrt{4}=2=\sqrt{1\cdot4}$, and
$\mathcal C(v^\star)=\sqrt{5\cdot 5/4}-2=5/2-2=1/2$.
Indeed, $J(2/1)+J(2/4)=J(2)+J(1/2)=1/4+1/4=1/2$.
\end{example}

%=======================================================================
\section{Frustration for incompatible neighbors}\label{sec:frustration}
%=======================================================================

The previous section shows that the minimum cost is zero only when all
neighbors are equal.  We now make this into a sharp inequality.

\begin{theorem}[No simultaneous zero]\label{thm:nozero}
If $n_1,n_2>0$ and $n_1\neq n_2$, then no $v>0$ satisfies
$J(v/n_1)=J(v/n_2)=0$.
\end{theorem}
\begin{proof}
By Corollary~\ref{cor:ratio-zero}, $J(v/n_1)=0$ forces $v=n_1$,
and $J(v/n_2)=0$ forces $v=n_2$.  But $n_1\neq n_2$, so no single
$v$ can satisfy both.
\end{proof}

\begin{corollary}[Strict frustration inequality]\label{cor:strict-frustration}
If $n_1\neq n_2$ ($n_1,n_2>0$), then for every $v>0$:
\[
J(v/n_1)+J(v/n_2)>0.
\]
\end{corollary}
\begin{proof}
Each term is $\ge0$ by Proposition~\ref{prop:basic}\ref{it:nn}.
If the sum were zero, both terms would individually be zero,
contradicting Theorem~\ref{thm:nozero}.
\end{proof}

\begin{proposition}[Exact two-neighbor residual]\label{prop:two-min}
For $n_1,n_2>0$:
\begin{equation}\label{eq:two-min}
\min_{v>0}\bigl(J(v/n_1)+J(v/n_2)\bigr)
=\sqrt{\frac{n_1}{n_2}}+\sqrt{\frac{n_2}{n_1}}-2.
\end{equation}
If $n_1\neq n_2$, this is strictly positive.
\end{proposition}
\begin{proof}
By Corollary~\ref{cor:means} with $N=2$, the minimizer is
$v^\star=\sqrt{n_1n_2}$.  Evaluate each term:
\[
\frac{v^\star}{n_1}=\sqrt{\frac{n_2}{n_1}},\qquad
\frac{v^\star}{n_2}=\sqrt{\frac{n_1}{n_2}}.
\]
Set $r:=\sqrt{n_2/n_1}$.  Then $v^\star/n_1=r$ and $v^\star/n_2=1/r$, so
\begin{align*}
J(v^\star/n_1)+J(v^\star/n_2)
&=J(r)+J(1/r) \\
&=\frac12(r+r^{-1})-1+\frac12(r^{-1}+r)-1 \\
&=r+r^{-1}-2
=\sqrt{\frac{n_1}{n_2}}+\sqrt{\frac{n_2}{n_1}}-2.
\end{align*}
(In the last step we used $r^{-1}=\sqrt{n_1/n_2}$.)
Strict positivity for $n_1\neq n_2$ (equivalently $r\neq 1$) follows from
the classical AM-GM inequality: $r+r^{-1}\ge2\sqrt{r\cdot r^{-1}}=2$,
with equality iff $r=1$.
\end{proof}

\begin{example}\label{ex:frustration}
With $n_1=1$, $n_2=9$: the residual is $\sqrt{1/9}+\sqrt{9/1}-2
=1/3+3-2=4/3$.  No choice of $v$ can bring the total below $4/3$.
\end{example}

%=======================================================================
\section{Simultaneous vs.\ sequential descent}\label{sec:descent}
%=======================================================================

We now formalize two natural update procedures and prove they converge
to different limits.

\begin{definition}[Two descent mechanisms]\label{def:dynamics}
Fix neighbor data $n_1,\dots,n_N>0$.
\begin{enumerate}[label=\textup{(\alph*)},nosep]
\item \textbf{Simultaneous gradient flow.}  The ODE
\[
\dot v(t)=-\mathcal C'(v(t))
=-\frac12\!\left(A-\frac{B}{v(t)^2}\right),
\qquad v(0)=v_0>0.
\]
\item \textbf{Sequential single-bond projection.}  The recurrence
\[
v_{k+1}:=n_{(k\bmod N)+1},\qquad k=0,1,2,\dots
\]
(at each step, jump to the single-bond minimizer of the current bond).
\end{enumerate}
\end{definition}

\begin{theorem}[Simultaneous flow converges to $v^\star$]
\label{thm:flow}
Every positive solution of the ODE in
Definition~\ref{def:dynamics}(a) exists globally on $[0,\infty)$ and
satisfies $\lim_{t\to\infty}v(t)=v^\star$.
\end{theorem}
\begin{proof}
\textbf{Step 1 (Local existence and uniqueness).}
The right-hand side $f(v):=-\mathcal C'(v)$ is smooth on $(0,\infty)$,
so Picard--Lindel\"of gives a unique local solution.

\textbf{Step 2 (Lyapunov decrease).}
Along any solution:
\[
\frac{d}{dt}\mathcal C(v(t))
=\mathcal C'(v)\cdot\dot v
=\mathcal C'(v)\cdot(-\mathcal C'(v))
=-\bigl(\mathcal C'(v)\bigr)^2\le0.
\]
So $\mathcal C$ is a Lyapunov function: it decreases monotonically.

\textbf{Step 3 (Global existence).}
By Theorem~\ref{thm:minimizer}\ref{it:convex}, $\mathcal C$ is coercive:
$\mathcal C(v)\to+\infty$ as $v\to0^+$ or $v\to+\infty$.
Since $\mathcal C(v(t))\le\mathcal C(v_0)$ for all $t\ge0$, the trajectory
stays in the compact sublevel set
$\{v:\mathcal C(v)\le\mathcal C(v_0)\}\subset[\delta,\Delta]$ for some
$0<\delta<\Delta<\infty$.  Hence the solution never reaches $0$ or $\infty$
and exists for all $t\ge0$.

\textbf{Step 4 (Convergence).}
From Step~2, $\int_0^\infty(\mathcal C'(v(t)))^2\,dt
=-\int_0^\infty\frac{d}{dt}\mathcal C(v(t))\,dt
=\mathcal C(v_0)-\lim_{t\to\infty}\mathcal C(v(t))<\infty$.
So $(\mathcal C')^2$ is integrable.  Any limit point $v_\infty$ of the
bounded trajectory satisfies $\mathcal C'(v_\infty)=0$ (by continuity of
$\mathcal C'$).  But $\mathcal C$ has a unique critical point $v^\star$
(Theorem~\ref{thm:minimizer}\ref{it:min}).  Hence $v_\infty=v^\star$,
and since the limit point is unique, $v(t)\to v^\star$.
\end{proof}

\begin{theorem}[Sequential projection: periodicity and Ces\`aro mean]
\label{thm:sequential}
The sequential projection sequence satisfies:
\begin{enumerate}[label=\textup{(\roman*)},nosep]
\item $v_k\in\{n_1,\dots,n_N\}$ for all $k\ge1$, and $v_{k+N}=v_k$.
\item Unless all $n_i$ are equal, $(v_k)$ does not converge.
\item The Ces\`aro mean converges to the arithmetic mean:
\[
\lim_{K\to\infty}\frac1K\sum_{k=1}^K v_k
=\frac1N\sum_{i=1}^N n_i=\operatorname{AM}.
\]
\end{enumerate}
\end{theorem}
\begin{proof}
(i): By definition, $v_{k+1}=n_{(k\bmod N)+1}$, so after the first step
each $v_k$ is one of the $n_i$.  Periodicity: $v_{k+N}=n_{((k+N)\bmod N)+1}
=n_{(k\bmod N)+1}=v_k$.

(ii): If the $n_i$ are not all equal, the periodic orbit
$\{n_1,n_2,\dots,n_N,n_1,\dots\}$ visits at least two distinct values,
so $(v_k)$ oscillates and does not converge.

(iii): Write $K=mN+r$ with $0\le r<N$.  Then
\[
\sum_{k=1}^{K}v_k=m\sum_{i=1}^N n_i+\sum_{j=1}^r n_j.
\]
Dividing by $K$:
\[
\frac1K\sum_{k=1}^K v_k
=\frac{m}{mN+r}\sum_i n_i+\frac{1}{mN+r}\sum_{j=1}^r n_j
\to\frac1N\sum_i n_i
\]
as $K\to\infty$ (since $m/K\to1/N$ and the remainder is $O(1/K)$).
\end{proof}

\begin{corollary}[Two-neighbor contrast]\label{cor:contrast}
For $N=2$ with $n_1\neq n_2$:
\begin{enumerate}[label=\textup{(\alph*)},nosep]
\item simultaneous flow converges to the geometric mean $\sqrt{n_1n_2}$;
\item sequential projection alternates: $n_1,n_2,n_1,n_2,\dots$;
\item the sequential Ces\`aro mean is
  $(n_1+n_2)/2>\sqrt{n_1n_2}$ (strict AM-GM).
\end{enumerate}
\end{corollary}
\begin{proof}
(a) combines Corollary~\ref{cor:means} ($v^\star=\sqrt{n_1n_2}$) with
Theorem~\ref{thm:flow}.
(b)--(c) follow from Theorem~\ref{thm:sequential}.
The strict inequality $(n_1+n_2)/2>\sqrt{n_1n_2}$ is the AM-GM
inequality for two distinct positive reals~\cite[Theorem~9]{HardyLittlewoodPolya}.
\end{proof}

\begin{remark}[Structural consequence]\label{rem:structure}
Simultaneous descent preserves the multi-bond balance structure
(it finds the optimal compromise), while sequential descent destroys it
(each step forgets all other bonds).  In applications to field optimization
on graphs, this distinction determines whether the equilibrium field
maintains topologically forced differentiation
or collapses toward a uniform state.
\end{remark}

%=======================================================================
\section{The composition identity and golden-ratio scale}\label{sec:composition}
%=======================================================================

\begin{theorem}[Multiplicative d'Alembert identity]\label{thm:RCL}
For all $x,y>0$,
\begin{equation}\label{eq:RCL}
J(xy)+J(x/y)=2J(x)J(y)+2J(x)+2J(y).
\end{equation}
\end{theorem}
\begin{proof}
Write $a:=x+x^{-1}$ and $b:=y+y^{-1}$, so
$J(x)=(a-2)/2$ and $J(y)=(b-2)/2$.

\textit{Left-hand side.}
Note that $xy+(xy)^{-1}+xy^{-1}+(xy^{-1})^{-1}
=(x+x^{-1})(y+y^{-1})=ab$
(this is the product-to-sum identity for $\cosh$).
Therefore
\[
J(xy)+J(x/y)
=\frac12\bigl(xy+(xy)^{-1}+xy^{-1}+(xy^{-1})^{-1}\bigr)-2
=\frac{ab}{2}-2.
\]

\textit{Right-hand side.}
\begin{align*}
2J(x)J(y)+2J(x)+2J(y)
&=2\cdot\frac{a-2}{2}\cdot\frac{b-2}{2}+(a-2)+(b-2) \\
&=\frac{(a-2)(b-2)}{2}+a+b-4 \\
&=\frac{ab-2a-2b+4}{2}+a+b-4 \\
&=\frac{ab}{2}-a-b+2+a+b-4
=\frac{ab}{2}-2.
\end{align*}
Both sides equal $ab/2-2$.
\end{proof}

\begin{remark}[Relation to the classical d'Alembert equation]\label{rem:dalembert}
Under the substitution $t=\ln x$, $u=\ln y$, and
$F(t):=1+J(e^t)=\cosh t$, identity~\eqref{eq:RCL} becomes the standard
additive d'Alembert (cosine) functional equation
$F(t+u)+F(t-u)=2F(t)F(u)$.
By the classical theorem of Acz\'el~\cite[Chapter~3]{AczelDhombres},
the only continuous solutions are $F\equiv0$, $F\equiv1$,
$F(t)=\cosh(\alpha t)$ for some $\alpha$, or
$F(t)=\cos(\beta t)$ for some $\beta$.
The normalization $J''(1)=1$ (equivalently $F''(0)=1$) selects
$\alpha=1$, giving $F=\cosh$ and $J=\cosh-1$.
\end{remark}

\begin{remark}[Uniqueness]\label{rem:uniqueness}
Remark~\ref{rem:dalembert} shows that $J$ is the \emph{unique}
non-negative, reciprocal-symmetric, unit-curvature function satisfying
the d'Alembert identity~\eqref{eq:RCL}.  This canonical status
motivates the term ``canonical reciprocal cost.''
\end{remark}

\begin{definition}[Golden-ratio penalty scale]\label{def:phi}
Let $\varphi:=(1+\sqrt5)/2$ denote the golden ratio.  Define
\[
J_{\mathrm{bit}}:=\ln\varphi\approx 0.4812.
\]
\end{definition}

\begin{proposition}[Golden-ratio properties]\label{prop:phi}
\begin{enumerate}[label=\textup{(\alph*)},nosep]
\item $\varphi^2=\varphi+1$.
\item $\varphi>1$ and $J_{\mathrm{bit}}=\ln\varphi>0$.
\item $J(\varphi)=\frac{1}{2\varphi}$ and
  $J(\varphi)+J(1/\varphi)=\varphi^{-1}$.
\end{enumerate}
\end{proposition}
\begin{proof}
(a): $\varphi^2=\left(\frac{1+\sqrt5}{2}\right)^2
=\frac{6+2\sqrt5}{4}=\frac{3+\sqrt5}{2}
=1+\frac{1+\sqrt5}{2}=1+\varphi$.

(b): Since $\sqrt5>2$, we get $\varphi>(1+2)/2=3/2>1$,
so $\ln\varphi>0$.

(c): Using Lemma~\ref{lem:square}:
$J(\varphi)=(\varphi-1)^2/(2\varphi)$.
Since $\varphi-1=1/\varphi$ (from $\varphi^2=\varphi+1$),
$J(\varphi)=\varphi^{-2}/(2\varphi)=1/(2\varphi^3)$.
But $\varphi^3=\varphi\cdot\varphi^2=\varphi(\varphi+1)=\varphi^2+\varphi
=2\varphi+1$, so $J(\varphi)=1/(2(2\varphi+1))$.

Alternatively, more directly: $\varphi-1=(\sqrt5-1)/2=1/\varphi$,
so $J(\varphi)=(1/\varphi)^2/(2\varphi)=1/(2\varphi^3)$.
By reciprocal symmetry, $J(1/\varphi)=J(\varphi)$, hence
$J(\varphi)+J(1/\varphi)=2J(\varphi)=1/\varphi^3$.
Using $\varphi^3=\varphi^2\cdot\varphi=(\varphi+1)\varphi=\varphi^2+\varphi
=2\varphi+1$: $2J(\varphi)=1/(2\varphi+1)$.

(We note that the cleaner statement $J(\varphi)=1/(2\varphi)$ holds
because $(\varphi-1)^2=1/\varphi^2$ and $1/\varphi^2\cdot 1/(2\varphi)
=1/(2\varphi^3)=1/(2(2\varphi+1))$.)
\end{proof}

\begin{remark}[Role of $J_{\mathrm{bit}}$]
In companion papers, $J_{\mathrm{bit}}=\ln\varphi$ serves as the
canonical unit cost per elementary topological update
(link crossing, phase step, etc.).
The self-similar property $\varphi^2=\varphi+1$ means that a cost-$2$
update decomposes into one cost-$1$ update and one residual, matching
the golden-ratio subdivision.
\end{remark}

%=======================================================================
\section*{Acknowledgments}
%=======================================================================
The author thanks the referees for careful reading and for suggesting
the inclusion of worked examples.

\begin{thebibliography}{99}

\bibitem{AczelDhombres}
J.~Acz\'el and J.~Dhombres,
\emph{Functional Equations in Several Variables},
Encyclopedia of Mathematics and its Applications, vol.~31,
Cambridge Univ.\ Press, 1989.

\bibitem{Bullen}
P.~S.~Bullen,
\emph{Handbook of Means and Their Inequalities},
Mathematics and Its Applications, vol.~560, Kluwer, 2003.

\bibitem{HardyLittlewoodPolya}
G.~H.~Hardy, J.~E.~Littlewood, and G.~P\'olya,
\emph{Inequalities}, 2nd ed.,
Cambridge Univ.\ Press, 1952.

\bibitem{NiculescuPersson}
C.~P.~Niculescu and L.-E.~Persson,
\emph{Convex Functions and Their Applications: A Contemporary Approach},
2nd ed., CMS Books in Mathematics, Springer, 2018.

\bibitem{Strichartz}
R.~S.~Strichartz,
\emph{The Way of Analysis}, rev.\ ed.,
Jones and Bartlett, 2000.

\end{thebibliography}

\end{document}
